\documentclass[a4paper,12pt]{article}
\usepackage[T2A]{fontenc}
\usepackage[utf8]{inputenc}
\usepackage[russian]{babel}
\usepackage{amsmath,amssymb,mathtools}
\usepackage{helvet}
\renewcommand{\familydefault}{\sfdefault}
\usepackage{icomma}
\usepackage[left=23mm,right=23mm,top=23mm,bottom=24mm]{geometry}
\usepackage{microtype}
\usepackage{xcolor}
\usepackage{tikz}
\usetikzlibrary{calc,arrows.meta,positioning,shapes.geometric}
\usepackage{tabularx,booktabs}
\usepackage{enumitem}
\usepackage{fancyhdr}
\usepackage{setspace}
\usepackage{parskip}
\setlength{\parskip}{4.5pt}
\usepackage{hyperref}
\hypersetup{hidelinks}
\definecolor{accent}{HTML}{177277}
\definecolor{darkaccent}{HTML}{194F55}
\definecolor{textgray}{HTML}{2E3537}
\color{textgray}
\setstretch{1.11}
\setlength{\headheight}{16pt}
\pagestyle{fancy}
\fancyhf{}
\fancyhead[L]{\small Ксения Васильченко}
\fancyhead[R]{\small 22.10.26}
\fancyfoot[C]{\small\thepage}
\renewcommand{\headrulewidth}{0.35pt}
\setlist[enumerate]{leftmargin=*,itemsep=.7em,topsep=.4em}
\newcommand{\sect}[1]{\section*{\color{darkaccent}#1}}
\newcommand{\learn}[1]{\par\noindent$\square$\enspace #1\par}
\newcommand{\remember}[1]{\par\smallskip\noindent\textcolor{darkaccent}{\textbf{Важно.}}\enspace #1\par\smallskip}
\newcommand{\mix}[3]{#1\,\frac{#2}{#3}}
\newcommand{\task}[1]{\item #1\par\vspace{2mm}}

\begin{document}
\hypersetup{pageanchor=false}
\begin{titlepage}
\thispagestyle{empty}
\centering
\vspace*{28mm}
{\color{accent}\large МАТЕМАТИКА \textbullet\ 6 КЛАСС\par}
\vspace{19mm}
{\color{darkaccent}\Huge\bfseries Чек-лист\par}
\vspace{9mm}
{\LARGE Деление дробей\par}
\vspace{4mm}
{\Large и уравнения с дробями\par}
\vspace{8mm}
{\large Обыкновенные, смешанные и десятичные дроби\par}
{\large Порядок действий \textbullet\ Неизвестный компонент\par}
\vspace{22mm}
{\large 22 октября 2026 года\par}
\vfill
{\large Лёвин Артём Александрович\par}
\vspace{3mm}
{\normalsize эксклюзивно для Ксении Васильченко\par}
\vspace{34mm}
\begin{tikzpicture}[remember picture,overlay]
\draw[accent!58,line width=1.25pt]
  ([xshift=25mm,yshift=20mm]current page.south west)
  .. controls ([xshift=79mm,yshift=42mm]current page.south west)
           and ([xshift=-77mm,yshift=4mm]current page.south east)
  .. ([xshift=-25mm,yshift=20mm]current page.south east);
\end{tikzpicture}
\end{titlepage}
\hypersetup{pageanchor=true}

\sect{1. Деление дробей: главное правило}
\learn{Вы заменяете деление умножением на \textbf{дробь, обратную делителю}.}
\learn{Вы проверяете, что делитель отличен от нуля.}

Если $b\ne0$, $c\ne0$, $d\ne0$, то
\[
\boxed{\frac ab:\frac cd=\frac ab\cdot\frac dc}.
\]
\textbf{Первая дробь остаётся прежней; переворачивается только вторая.}

\textbf{Пример 1.} Разделим две обыкновенные дроби:
\[
\frac13:\frac25=\frac13\cdot\frac52=\boxed{\frac56}.
\]
\textbf{Пример 2.} После замены деления умножением можно сокращать:
\[
\frac59:\frac49=\frac59\cdot\frac94=\boxed{\frac54=\mix{1}{1}{4}}.
\]
\remember{Делить на ноль нельзя. Например, выражение $\frac25:0$ не имеет значения.}

\sect{2. Деление с натуральным числом}
\learn{Вы записываете натуральное число как дробь со знаменателем~$1$.}
\textbf{Дробь делим на число:}
\[
\frac13:2=\frac13:\frac21=\frac13\cdot\frac12=\boxed{\frac16}.
\]
\textbf{Число делим на дробь:}
\[
2:\frac13=\frac21:\frac13=\frac21\cdot\frac31=\boxed6.
\]

\sect{3. Деление смешанных чисел}
\learn{Вы сначала переводите смешанные числа в неправильные дроби.}
\[
\mix{a}{b}{c}=\frac{ac+b}{c}\qquad (a\ge0,\;0\le b<c).
\]
\textbf{Пример.} Найдём частное:
\[
\mix{2}{1}{5}:\frac13=\frac{11}{5}\cdot3
=\frac{33}{5}=\boxed{\mix{6}{3}{5}}.
\]
\remember{Сокращение делайте до умножения, когда это возможно. В конце сократите ответ и при необходимости выделите целую часть.}

\sect{4. Деление десятичной дроби на обыкновенную}
\learn{Вы переводите десятичную дробь в обыкновенную и сокращаете её.}
\learn{Вы правильно определяете знаменатель по числу знаков после запятой.}

\[
0{,}75:\frac14=\frac{75}{100}:\frac14
=\frac34\cdot4=\boxed3.
\]
\textbf{Полезно узнавать эти дроби сразу:}
\renewcommand{\arraystretch}{1.35}
\begin{center}
\begin{tabular}{@{}cccc@{}}
\toprule
$0{,}5=\dfrac12$ & $0{,}25=\dfrac14$ & $0{,}75=\dfrac34$ & $0{,}2=\dfrac15$\\
$0{,}8=\dfrac45$ & $0{,}125=\dfrac18$ & $0{,}08=\dfrac2{25}$ & $0{,}05=\dfrac1{20}$\\
\bottomrule
\end{tabular}
\end{center}
\textbf{Сравните:}
\[
0{,}8=\frac8{10}=\frac45,
\qquad
0{,}08=\frac8{100}=\frac2{25}.
\]
\remember{Один знак после запятой --- десятые; два знака --- сотые. Это существенно меняет ответ.}

\sect{5. Порядок действий}
\learn{Вы сначала выполняете действия в скобках.}
\learn{Вы записываете длинное выражение по действиям и сокращаете до умножения.}

\textbf{Пример.} Вычислим:
\[
\left(\mix{2}{1}{4}:\mix{3}{6}{7}\right)\cdot\mix{2}{1}{7}.
\]
\textbf{Первое действие (в скобках):}
\[
\frac94:\frac{27}{7}=\frac94\cdot\frac7{27}=\frac7{12}.
\]
\textbf{Второе действие:}
\[
\frac7{12}\cdot\frac{15}{7}
=\frac{15}{12}=\boxed{\frac54=\mix{1}{1}{4}}.
\]
\remember{В примере $\frac94:\frac{27}{7}$ знак деления исчезает только вместе с переворачиванием делителя.}

\sect{6. Как находить неизвестный компонент}
\learn{Вы определяете, какое выражение является неизвестным слагаемым, множителем или вычитаемым.}
\learn{Вы сначала упрощаете известную часть уравнения.}

\renewcommand{\arraystretch}{1.4}
\begin{center}
\begin{tabularx}{\linewidth}{@{}>{\raggedright\arraybackslash}X>{\raggedright\arraybackslash}X@{}}
\toprule
\textbf{Вид уравнения} & \textbf{Как найти неизвестное}\\
\midrule
$a\cdot x=c$, $a\ne0$ & $x=c:a$\\
$x+b=c$ & $x=c-b$\\
$a-x=c$ & $x=a-c$\\
$x-b=c$ & $x=c+b$\\
\bottomrule
\end{tabularx}
\end{center}

\sect{7. Уравнения: решайте по этапам}
\textbf{Пример 1. Неизвестный множитель:}
\[
\frac27z=\mix{1}{1}{7},\qquad
z=\frac87:\frac27=\frac87\cdot\frac72=\boxed4.
\]

\textbf{Пример 2. Сначала сложение известных дробей:}
\begin{align*}
\frac59m&=\frac5{18}+\frac12,\\
\frac59m&=\frac5{18}+\frac9{18}=\frac79,\\
m&=\frac79:\frac59=\boxed{\frac75}.
\end{align*}

\textbf{Пример 3. Два этапа поиска неизвестного:}
\begin{align*}
\frac49b+\frac37&=1,\\
\frac49b&=1-\frac37=\frac47,\\
b&=\frac47:\frac49=\boxed{\frac97}.
\end{align*}

\textbf{Пример 4. Неизвестное вычитаемое:}
\begin{align*}
\frac12-\frac37x&=\frac14,\\
\frac37x&=\frac12-\frac14=\frac14,\\
x&=\frac14:\frac37=\boxed{\frac7{12}}.
\end{align*}

\textbf{Пример 5. Если неизвестное выражение стоит первым:}
\begin{align*}
\frac14x-2&=\frac35,\\
\frac14x&=\frac35+2=\frac{13}{5},\\
x&=\frac{13}{5}:\frac14
=\frac{13}{5}\cdot4=\boxed{\frac{52}{5}}.
\end{align*}
\remember{Следите за порядком деления: для $a\cdot x=c$ нужно вычислить $c:a$. Переписав уравнение, используйте \emph{новую} правую часть, а не число из первоначального условия.}

\sect{8. Самопроверка перед ответом}
\learn{Записан ли обратный \emph{делителю} множитель?}
\learn{Верно ли переведены смешанные и десятичные дроби?}
\learn{Сначала выполнены действия в скобках и вычислена известная часть?}
\learn{После решения уравнения подстановка даёт верное равенство?}

\clearpage
\thispagestyle{plain}
\begin{center}
{\color{darkaccent}\LARGE\bfseries Алгоритм: деление дробей}
\end{center}
\vspace{9mm}
\begin{center}
\begin{tikzpicture}[
  node distance=13mm,
  >={Stealth[length=2.7mm]},
  every node/.style={font=\small,align=center},
  terminator/.style={draw=darkaccent,rectangle,rounded corners=5pt,minimum width=51mm,minimum height=11mm,fill=accent!5,thick},
  action/.style={draw=accent,rounded corners=2mm,text width=87mm,minimum height=14mm,inner sep=3mm,thick},
  choice/.style={draw=darkaccent,diamond,aspect=3.8,text width=48mm,inner sep=2mm,thick},
  side/.style={draw=accent,rounded corners=2mm,text width=30mm,minimum height=14mm,inner sep=3mm,thick},
  path/.style={->,line width=.95pt,darkaccent}
]
\node[terminator] (start) {Начало};
\node[choice,below=of start] (zero) {Делитель равен нулю?};
\node[side,right=19mm of zero] (stop) {Деление\ невозможно};
\node[action,below=15mm of zero] (convert) {Приведите числа к обыкновенным дробям:\ смешанные --- к неправильным, целые --- к $n/1$};
\node[action,below=of convert] (reciprocal) {Замените деление умножением\ на дробь, обратную делителю};
\node[action,below=of reciprocal] (reduce) {Сократите общие множители числителя\ и знаменателя, затем перемножьте};
\node[action,below=of reduce] (format) {Сократите результат; если нужно,\ выделите целую часть};
\node[terminator,below=of format] (finish) {Ответ};
\draw[path] (start)--(zero);
\draw[path] (zero)--node[left,pos=.5]{нет}(convert);
\draw[path] (zero)--node[above,pos=.54]{да}(stop);
\draw[path] (convert)--(reciprocal);
\draw[path] (reciprocal)--(reduce);
\draw[path] (reduce)--(format);
\draw[path] (format)--(finish);
\end{tikzpicture}
\end{center}

\clearpage
\sect{Тренировочный блок --- Путь А}
\textbf{Выполните 10 заданий.} Работайте по нарастающей сложности.
Для вычислений указывайте преобразование деления в умножение;
для уравнений записывайте промежуточные этапы.
\begin{enumerate}
\task{Вычислите: $\displaystyle\frac25:\frac34$.}
\task{Вычислите: $\displaystyle\frac37:6$.}
\task{Вычислите: $\displaystyle4:\frac23$.}
\task{Вычислите: $\displaystyle\mix{1}{1}{2}:\frac34$.}
\task{Вычислите: $\displaystyle0{,}75:\frac38$.}
\task{Вычислите: $\displaystyle0{,}08:\frac25$. Объясните, почему знаменатель при переводе десятичной дроби равен $100$.}
\task{Вычислите по действиям:
$\displaystyle\left(\mix{2}{1}{4}:\mix{3}{6}{7}\right)\cdot\mix{2}{1}{7}$.}
\task{Решите уравнение: $\displaystyle\frac27z=\mix{1}{1}{7}$.}
\task{Решите уравнение: $\displaystyle\frac49b+\frac37=1$.}
\task{Решите уравнение: $\displaystyle\frac14x-2=\frac35$. Выполните проверку подстановкой.}
\end{enumerate}

\clearpage
\sect{Ответы к тренировочному блоку}
\vspace{5mm}
\renewcommand{\arraystretch}{1.8}
\begin{center}
\begin{tabularx}{\linewidth}{@{}c>{\raggedright\arraybackslash}X@{}}
\toprule
\textbf{№} & \textbf{Ответ}\\
\midrule
1 & $\dfrac8{15}$\\
2 & $\dfrac1{14}$\\
3 & $6$\\
4 & $2$\\
5 & $2$\\
6 & $\dfrac15$\\
7 & $\dfrac54=\mix{1}{1}{4}$\\
8 & $z=4$\\
9 & $b=\dfrac97=\mix{1}{2}{7}$\\
10 & $x=\dfrac{52}{5}=\mix{10}{2}{5}$\\
\bottomrule
\end{tabularx}
\end{center}
\vspace{8mm}
\remember{В уравнениях проверьте найденное число подстановкой. При делении
на положительную дробь меньше $1$ результат больше делимого;
при делении на число больше $1$ --- меньше.}
\end{document}
